\subsection{A process dies while its peer is present}
\label{sec:ha:enabled:process_death}

A component dies on a machine that is otherwise healthy, and the launcher
restarts it, while a second instance of the same component runs elsewhere. Both
configurations apply to this case.

\subsubsection{What is at risk}
\label{sec:ha:enabled:process_death:risk}

The venue's ability to trade is not at risk, because the peer can act and trading
continues whatever the venue decides. What is at risk is \textbf{the pair}.

If the peer promotes, one instance serves the venue until somebody restores the
other, and the venue is exposed to a second failure for as long as that takes. If
the launcher restarts the dead instance instead, the pair is whole again in
seconds and the venue loses nothing.

Nothing the venue does can restore a machine that has died, so promoting the peer
is the only remedy for that failure. The venue can restore a process that has
died on a living machine, so promoting the peer in that case gives up a recovery
that was available. The venue must therefore treat the two failures differently.

\subsubsection{What the venue must do}
\label{sec:ha:enabled:process_death:behaviour}

\begin{requirement}{R-0080}{A process death on a healthy machine does not cost the pair}
  Where a component dies and is restarted on a machine that has not failed, the
  venue shall end with two instances running, as it began.
  \rationale{The measure of this requirement is whether the pair is whole again, rather than
  whether the venue kept trading. A promotion keeps the venue trading on its own,
  and leaves one instance down with nothing watching for a second failure.}
  \uncovered{no scenario asserts the pair is whole afterwards, only that trading continued}
\end{requirement}

\begin{requirement}{R-0081}{The peer does not promote while a restart could still succeed}
  An instance that has lost contact with the acting instance shall wait, before
  seeking to act, for at least as long as a supervised restart takes.
  \rationale{A shorter wait fails the venue over for a failure that did not require it. A much
  longer wait tolerates a genuinely dead machine for longer than it should. The
  period is therefore the time a restart takes plus a margin. That time is not
  known until restarts have been timed, so the period is measured rather than
  chosen, and it must be measured again whenever the components change.
  The peer cannot take over until its own promise to the acting instance and the
  arbiter's have both run out, which is at least one lease period after the last
  renewal. A restarted instance that kept its promises on disk asks again as soon
  as it starts, so a restart that completes within the lease period keeps the
  lead.}
  \covers{ha_test:26, ha_test:27, ha_test:28}
\end{requirement}

\begin{requirement}{R-0082}{An instance that keeps failing does not disturb a healthy peer}
  A restarted instance shall not take the entitlement to act from a peer that
  holds it and is present.
  \rationale{Without this requirement, an instance that fails repeatedly takes leadership on
  each restart, dies, and hands leadership back, interrupting a peer that was
  working correctly once per failure. \reqref{R-0062} states the arbiter's side of
  this rule, and this requirement states the instance's side.}
  \covers{ha_test:29, ha_test:30, ha_test:31}
\end{requirement}

\begin{requirement}{R-0125}{An instance that interrupts the venue twice is not restarted again}
  Where restarting an instance has already cost the venue one interruption, a
  second interruption from that instance shall leave it stopped rather than
  restarted.
  \rationale{An interruption is a period in which a member can neither place an order nor
  withdraw one, and is told nothing while it lasts. The venue can account for one
  interruption: a component failed, and recovery took as long as it took. A second
  interruption from the same instance shows that the instance is unfit to serve,
  and restarting it again produces a third.

  Stopping the instance is what allows the pair to do its work. An instance that
  the supervisor keeps restarting returns before the arbiter can grant the
  entitlement to its peer, so the peer never acts. The venue then takes an
  unbounded series of interruptions instead of the single failover the pair was
  built for. Leaving the failing instance stopped produces the outcome the member
  should have had at the first fault.

  \textbf{The count is taken over a session rather than over a period.} A time
  window cannot be explained to the member it affects, because it makes two
  failures count together when they fall close in time and separately when they do
  not. The requirement is stated of interruptions rather than of failed starts for a
  related reason: an instance that dies every few minutes starts correctly every
  time.}
  \uncovered{no scenario flaps an instance and asserts it is left stopped}
\end{requirement}

\begin{requirement}{R-0126}{The venue does not run with no instance of a component without saying so}
  Where the last instance of a component is stopped and not restarted, the venue
  shall report that it cannot perform that component's work, rather than continuing
  in a state where nothing does it.
  \rationale{Both instances may share the cause of the failure: an input that kills whatever
  reads it, a build, or a configuration both instances load. The peer then meets
  that cause in its turn and is stopped in its turn. The venue has no instance of
  that component left, and nothing is trying to produce one.

  This failure looks like health. Every process that remains is running normally,
  no alert names a component, and only somebody who counts the processes can see
  what is missing. Reporting the condition costs an announcement. Not reporting it
  costs the time between the last instance stopping and somebody noticing, and that
  time is unbounded.}
  \uncovered{no scenario stops every instance of a component and asserts the venue says so}
\end{requirement}

\subsubsection{When the restart and the promotion race}
\label{sec:ha:enabled:process_death:race}

The two remedies run at the same time, and neither detects the other. The
launcher restarts the instance that died while the peer counts down to asking
whether it may act. \textbf{Exactly one instance ends up acting, whichever of the
two happens first.}

\textbf{The restart completes first.} The returning instance holds no position
and asks which position it holds. No instance holds the entitlement, so the
arbiter grants it to the returning instance, which resumes acting. The peer is
still waiting, finds the entitlement held, and remains the follower. The
entitlement did not move, and the pair is whole.

\textbf{The peer asks first.} The arbiter grants the peer the entitlement, and
the peer begins acting. The returning instance asks, finds the entitlement held
by an instance that is present, and becomes the follower. Neither its configured
name nor its having held the entitlement minutes earlier affects that answer.

\textbf{Both ask at once.} One party answers both claims. The arbiter grants the
entitlement to one instance, records what it granted, and refuses the other claim
against that record.

In every ordering, exactly one instance ends up acting, and no instance has to
detect what the other was doing. That is what makes the outcome certain rather
than merely likely.

\begin{requirement}{R-0090}{A returning instance assumes nothing and asks}
  An instance that has restarted shall hold no position until it is told which
  it holds, whatever position it held before and whatever it is configured as.
  \rationale{An instance that assumes it leads, because it led before or because of its
  configured name, becomes a second acting instance whenever the peer has already
  been promoted. Nothing has to go wrong with the peer for that to happen.}
  \covers{ha_test:24, ha_test:25}
\end{requirement}

\begin{requirement}{R-0091}{A returning instance that finds the position vacant takes it}
  Where an instance returns and no instance holds the entitlement, it shall take
  it rather than wait.
  \rationale{A returning instance holds no position and asks, which is what stops a restart
  producing a second acting instance. This requirement states the other half of
  that rule. Where no instance holds the entitlement, the returning instance must
  take it. Without this, a restart that completed before its peer's promotion would
  leave both instances following and none acting.}
  \covers{ha_test:26, ha_test:27, ha_test:28}
\end{requirement}

\begin{requirement}{R-0092}{Simultaneous claims are resolved by one party, not by the claimants}
  Where a returning instance and its peer both seek the entitlement at once, at
  most one shall be granted it, and neither shall be required to have detected
  the other.
  \rationale{Two instances that must detect each other in order to stay out of each other's
  way will eventually fail to do so, and they will fail at the moment they are
  least able to communicate. Referring both claims to one party removes the need to
  detect anything.}
  \uncovered{no scenario times a restart to collide with a promotion}
\end{requirement}
